% Generated draft: requires mathematical and source review.
\documentclass[11pt]{article}
\usepackage[margin=25mm]{geometry}
\usepackage{fontspec}
\setmainfont{Latin Modern Roman}
\usepackage{amsmath,amssymb,mathtools,unicode-math}
\setmathfont{Latin Modern Math}
\usepackage{enumitem,needspace,xurl,hyperref}
\setlength{\parindent}{0pt}
\setlength{\parskip}{5pt}
\setlength{\emergencystretch}{3em}
\newcommand{\N}{\mathbb{N}}
\newcommand{\Z}{\mathbb{Z}}
\newcommand{\Q}{\mathbb{Q}}
\newcommand{\R}{\mathbb{R}}
\newcommand{\C}{\mathbb{C}}
\newcommand{\F}{\mathbb{F}}
\newcommand{\A}{\mathbb{A}}
\newcommand{\PP}{\mathbb{P}}
\newcommand{\E}{\mathbb{E}}
\newcommand{\eps}{\varepsilon}
\newcommand{\dd}{\,\mathrm{d}}
\begin{document}
\section*{UnsolvedMath ID: 800649}
\textbf{Source problem code:} catalogue entry
\subsection{Definitions and mathematical statement}
For a topological space \(X\), a family \(\mathcal U\) of open subsets of \(X\) is called large if
\[
\forall x\in X,\qquad \bigl|\{U\in\mathcal U:x\in U\}\bigr|\text{ is infinite}.
\]
Write \(\Lambda(X)\) for the source's intended class of large open covers of \(X\). The supplied excerpt does not specify whether these covers must be countable, whether \(X\) itself is excluded as a cover member, or which class of spaces is under consideration. These conventions must be recovered from the source before the formulation is finalized.

Relative to those conventions, \(X\) satisfies \(\mathsf{Split}(\Lambda,\Lambda)\) if
\[
\forall\mathcal U\in\Lambda(X)\ \exists\mathcal U_0,\mathcal U_1\subseteq\mathcal U,
\qquad
\mathcal U_0\cap\mathcal U_1=\varnothing
\quad\text{and}\quad
\mathcal U_0,\mathcal U_1\in\Lambda(X).
\]
Thus every eligible large open cover contains two disjoint subfamilies that are themselves eligible large open covers.

The question is whether \(\mathsf{Split}(\Lambda,\Lambda)\) is additive. Under the usual interpretation as preservation under finite unions of subspaces, this asks whether, for every admissible common ambient space \(Z\) and all admissible subspaces \(X,Y\subseteq Z\),
\[
\bigl[X\models\mathsf{Split}(\Lambda,\Lambda)
\ \land\ Y\models\mathsf{Split}(\Lambda,\Lambda)\bigr]
\ \Longrightarrow\
X\cup Y\models\mathsf{Split}(\Lambda,\Lambda),
\]
where \(X\), \(Y\), and \(X\cup Y\) carry their subspace topologies. This is a conditional reconstruction: the excerpt does not identify the admissible spaces or define ``additive,'' so it does not establish whether finite unions, countable unions, or another source-specific convention is intended.

The source presents this among the remaining preservation questions for splittability properties under set-theoretic operations and says that the accompanying problems are settled for properties omitted from them. That statement is historical source context, not an independently verified current classification.
\subsection{Short English statement}
Does the union of two spaces satisfying \(\mathsf{Split}(\Lambda,\Lambda)\) again satisfy it? Here splitting means finding two disjoint large-cover subfamilies inside each eligible large open cover. This finite-union reading is provisional: the intended spaces, cover conventions, and meaning of additivity require confirmation from the source. The registry labels the problem open; that label has not been independently verified.
\subsection{Sources}
\begin{itemize}
\item Local UnsolvedMath registry, permanent ID \(800649\), source problem code \texttt{catalogue entry}, title ``Selection Principles Milestone Problems --- Problem 22 ([splittability, Problem 6.8]).'' Supplied upstream revision: \texttt{d371ee917ac9647fbb06cd6fc529deec6f3dda82}. The registry attributes the problem to Tsaban, with proposed year \(2003\), and supplies status ``open.''
\item Supplied source list: Tsaban, \emph{Selection Principles in Mathematics: A Milestone of Open Problems} (2003), Problem 22. Supplied URL: \url{https://arxiv.org/abs/math/0312182}. The record reports access on \texttt{2026-10-16} and extraction by \texttt{arxiv\_latex}; these are supplied provenance details, not an assertion of independent access or verification.
\item Supplied attribution/context: ``[splittability, Problem 6.8].'' The supplied record does not resolve the reference label ``splittability'' into a bibliographic entry.
\end{itemize}
\end{document}
